\documentclass[11pt,a4paper]{article}
\usepackage[margin=1in]{geometry}
\usepackage{amsmath,amssymb,amsthm}
\usepackage{booktabs}
\usepackage{hyperref}
\usepackage{xcolor}

\newtheorem{conjecture}{Conjecture}
\newtheorem{theorem}{Theorem}
\newtheorem{proposition}{Proposition}
\newtheorem{lemma}{Lemma}
\newtheorem{definition}{Definition}

\DeclareMathOperator{\ord}{ord}

\title{Furstenberg's $\times 2, \times 3$ Conjecture:\\
A Computational-Diophantine Approach via Baker Separation\\
and Partition Entropy}
\author{John A. Janik\\
john.janik@gmail.com}
\date{February 2026}

\begin{document}
\maketitle

\begin{abstract}
We present large-scale computational evidence and structural analysis
bearing on Furstenberg's 1967 conjecture that the only Borel probability
measures on $\mathbb{T} = \mathbb{R}/\mathbb{Z}$ jointly invariant under
$T_2\colon x \mapsto 2x$ and $T_3\colon x \mapsto 3x$ are Lebesgue
measure and measures supported on finite orbits.
Our approach transfers the Baker--cellError machinery developed for
the Collatz conjecture to the Furstenberg setting, producing three
categories of results.

\emph{Phase~1 (Orbit structure):}
We scan $4.55 \times 10^8$ primes up to $10^{10}$, verifying that the
density of primes~$p$ for which $\langle 2, 3 \rangle = (\mathbb{Z}/p\mathbb{Z})^*$
converges to the Euler product
$\prod_\ell (1 - 1/(\ell^2(\ell-1))) = 0.69750\ldots$
to five significant figures.
We also compute continued-fraction convergents of $\log_2 3$
and multi-scale cellError resonant fractions, confirming that the
resonant count stays $O(1)$ at every scale (Baker separation).

\emph{Phase~2 (Entropy bridge):}
We compute the joint partition $\mathcal{P}_2^{(j)} \vee \mathcal{P}_3^{(k)}$
for $j \leq 24$, $k \leq 15$, finding that:
(i)~the entropy deficit $\log(\text{\#atoms}) - H$ is
\emph{bounded} at $\sim 0.18$ nats across all scales;
(ii)~the minimum atom width equals the B\'ezout floor $1/(2^j \cdot 3^k)$ exactly;
(iii)~a fundamental asymmetry holds: $\ord_{2^j}(3) = 2^{j-2}$ (index~2)
while $\ord_{3^k}(2) = \varphi(3^k)$ (index~1, primitive root), both
verified computationally to exponent~40.
These results constrain the entropy bridge argument: $T_3$-invariance
forces approximate uniformity on $\mathcal{P}_2^{(j)}$ atoms, which
together with Rudolph's theorem would close the conjecture.

\emph{Phase~2b (Spectral gap):}
We prove analytically and verify numerically that the transfer matrix
of $T_3$ on $\mathcal{P}_2^{(j)}$ has \emph{all} non-trivial eigenvalues
with modulus exactly~$1/3$, giving spectral gap $2/3$ \emph{independent of~$j$}.
The proof uses a telescoping product identity:
$\prod_{i=0}^{L-1} |1 + 2\cos\theta_i| = 1$ for every orbit of
$\times 3$ on $\mathbb{Z}/2^j\mathbb{Z}$.
This is verified to $j = 30$ via Fourier orbit products ($2^{30}$ elements)
and to $j = 10$ via direct matrix multiplication ($1024 \times 1024$ matrices,
4 orbit cycles each).
The uniform spectral gap gives exponential mixing:
$\|T_3^n p - \text{uniform}\|_2 \leq (1/3)^n \cdot \|p - \text{uniform}\|_2$.

\emph{Phase~2c (Joint constraints):}
We build the combined $T_2 + T_3$ constraint matrix across 2--3 consecutive
dyadic scales (up to $9216 \times 9216$ at $j = 10$) and compute exact
rank formulas via Gaussian elimination:
$\operatorname{rank} = 7N - 1$ (2-scale) or $17N - 1$ (3-scale),
with $T_2$-invariance killing exactly $N_{\text{finest}}$ degrees of freedom.
However, the atom-mass DOF at the coarsest scale remains $N/2 + 1$, showing
that the linear constraint system alone does not force $\mathbf{T}p = p$.
This identifies \emph{positivity} ($w_{k,s} \ge 0$) as the essential
non-linear ingredient for the entropy bridge, consistent with the structure
of Rudolph's theorem.
\end{abstract}


%% ====================================================================
\section{Introduction}
%% ====================================================================

Furstenberg's conjecture \cite{furstenberg1967} is one of the central
open problems in ergodic number theory.

\begin{conjecture}[Furstenberg, 1967]
The only Borel probability measures on $\mathbb{T} = \mathbb{R}/\mathbb{Z}$
simultaneously invariant under $T_2\colon x \mapsto 2x \pmod{1}$ and
$T_3\colon x \mapsto 3x \pmod{1}$ are Lebesgue measure~$\lambda$ and
measures supported on finite orbits (rationals with denominator coprime to~6).
\end{conjecture}

The state of the art includes:
\begin{itemize}
\item \textbf{Rudolph (1990):} If $\mu$ is ergodic, jointly $T_2$- and
  $T_3$-invariant, and $h_\mu(T_2) > 0$, then $\mu = \lambda$.
\item \textbf{Host (1995):} Simplified Rudolph's proof.
\item \textbf{Einsiedler--Katok--Lindenstrauss (2006):} Any
  counterexample has Hausdorff dimension~0.
\item \textbf{The zero-entropy case remains completely open.}
\end{itemize}

Our approach uses Baker's theorem on linear forms in logarithms---the
same effective Diophantine tool that drives the analysis of the Collatz
conjecture via ``cellError'' separation---to obtain quantitative
constraints on the joint partition structure of $T_2$ and $T_3$.
The goal is to show that Baker separation forces positive entropy for
any non-atomic jointly invariant measure, reducing Furstenberg to Rudolph.


%% ====================================================================
\section{Mathematical Framework}
%% ====================================================================

\subsection{The Fourier picture}

Characters of $\mathbb{T}$ are $\chi_n(x) = e^{2\pi i n x}$ for
$n \in \mathbb{Z}$.
If $\mu$ is $T_2$-invariant, then $\hat\mu(n) = \hat\mu(2n)$;
if also $T_3$-invariant, then $\hat\mu(n) = \hat\mu(3n)$.
Thus $\hat\mu$ is constant on orbits of $\langle \times 2, \times 3 \rangle$
acting on $\mathbb{Z}$.
For $n \neq 0$, the orbit $\mathcal{O}(n) = \{2^a 3^b n : a, b \geq 0\}$
is infinite by multiplicative independence of~2 and~3.
The orbit tends to $\infty$, so by Riemann--Lebesgue
$\hat\mu(n) = 0$ if $\mu$ is absolutely continuous.
For singular~$\mu$, Riemann--Lebesgue fails, and the question is
whether the Fourier coefficients can consistently remain nonzero
along the rapidly diverging orbit.

\subsection{Baker separation on the torus}

Baker's theorem \cite{baker1966} gives a lower bound on linear forms
in logarithms.  In the notation of \cite{rhin1987}:
\[
|a \log 2 + b \log 3| > \frac{C}{\max(|a|, |b|)^\kappa}
\quad\text{for } (a,b) \neq (0,0),
\]
with $\kappa = 4.125$ (Rhin's effective value).
Define the \emph{cellError}:
\[
\operatorname{cellError}(a, b) = a - \log_2 3 \cdot b.
\]
Baker gives $|\operatorname{cellError}(a,b)| > C / \max(|a|,|b|)^\kappa$,
ensuring that the 2-adic and 3-adic structures are always ``transverse.''

On the $(N \times N)$ torus, a cell $(a,b)$ is \emph{resonant} if
$|\operatorname{cellError}(a,b)| < 1/N$.
Baker bounds the resonant count at $O(1)$ (independent of~$N$),
so the resonant \emph{fraction} decays as $O(1/N)$.
This is the ``safe cell density'' of the Collatz framework,
now applied to the pure $\times 2, \times 3$ action.

\subsection{The entropy bridge conjecture}

\begin{conjecture}[Entropy Production]
\label{conj:entropy}
If $\mu$ is a non-atomic, ergodic Borel probability measure on $\mathbb{T}$
that is simultaneously $T_2$- and $T_3$-invariant, then
\[
h_\mu(T_2) = \lim_{j \to \infty} \frac{1}{j}\, H_\mu\!\!\left(
\bigvee_{i=0}^{j-1} T_2^{-i}(\mathcal{P}_2)\right) > 0,
\]
where $\mathcal{P}_2 = \{[0, \tfrac12), [\tfrac12, 1)\}$.
\end{conjecture}

If Conjecture~\ref{conj:entropy} holds, then Rudolph's theorem
immediately implies $\mu = \lambda$, closing Furstenberg's conjecture.

The partition $\mathcal{P}_2^{(j)} = \bigvee_{i=0}^{j-1} T_2^{-i}(\mathcal{P}_2)$
has atoms $\{[m/2^j, (m+1)/2^j) : 0 \leq m < 2^j\}$ with boundary
points at the 2-adic rationals $\{m/2^j\}$.
Similarly, $\mathcal{P}_3^{(k)}$ has boundaries at $\{m'/3^k\}$.
The joint partition $\mathcal{P}_2^{(j)} \vee \mathcal{P}_3^{(k)}$ overlays
both sets of boundaries.
Baker's theorem ensures that 2-adic and 3-adic boundaries never coincide
(except at~0), and quantifies their minimal separation.


%% ====================================================================
\section{Computational Methods}
%% ====================================================================

All computations were performed on an Intel Core Ultra 9~275HX
(24~cores, no hyperthreading) with 128\,GB DDR5, using OpenMP
parallelization.

\subsection{Phase~1: Prime orbit scan (\texttt{furstenberg\_orbits.c})}

For each prime $p$ with $\gcd(p, 6) = 1$ up to $P_{\max}$:
\begin{enumerate}
\item Factor $p - 1$ by trial division using precomputed primes to $10^5$.
\item Compute $\ord_p(2)$ and $\ord_p(3)$ via the factorization of $p-1$:
  start with $\text{ord} = p-1$, then for each prime factor $q$ of $p-1$,
  divide out~$q$ while $x^{\text{ord}/q} \equiv 1 \pmod{p}$.
\item Compute $|\langle 2, 3 \rangle| = \text{lcm}(\ord_p(2), \ord_p(3))$
  using the cyclic group identity.
\item Test full generation: $\langle 2, 3 \rangle = (\mathbb{Z}/p\mathbb{Z})^*$
  iff $\text{lcm} = p - 1$.
\end{enumerate}
The sieve and prime collection use bit-packed Eratosthenes (625\,MB at
$10^{10}$).
Checkpoints are saved every 200K primes via atomic rename.
At $10^{10}$, the program processes $7.9 \times 10^6$~primes/s
on 24~threads.

\subsection{Phase~2: Entropy computation (\texttt{furstenberg\_entropy.c})}

For each $(j, k)$ pair, the joint partition
$\mathcal{P}_2^{(j)} \vee \mathcal{P}_3^{(k)}$ has boundary points
$\{m/2^j\} \cup \{m'/3^k\}$.
In the common denominator $D = 2^j \cdot 3^k$, these become integers:
\[
\{m \cdot 3^k : 0 \leq m < 2^j\} \cup \{m' \cdot 2^j : 0 \leq m' < 3^k\}.
\]
We sort these $2^j + 3^k - 1$ integers (only~0 is shared), compute the
gaps between consecutive values, and derive:
\begin{itemize}
\item Atom count: $N_{\text{atoms}} = 2^j + 3^k - 1$.
\item Minimum and maximum gap in integer units.
\item Entropy: $H = -\sum_i w_i \ln w_i$ where $w_i = \text{gap}_i / D$.
\item Entropy deficit: $\ln(N_{\text{atoms}}) - H$.
\end{itemize}
For $j \leq 24$, $k \leq 15$, the boundary array fits in $\sim 250$\,MB.

Multiplicative orders $\ord_{2^j}(3)$ and $\ord_{3^k}(2)$ are computed
by factoring $\varphi(2^j)$ and $\varphi(3^k)$ and iteratively dividing
out prime factors while the power test passes, using 128-bit modular
arithmetic.

\subsection{Phase~2b: Spectral verification (\texttt{furstenberg\_spectrum.c})}

The transfer matrix spectrum is verified in two ways:
(1)~Fourier orbit products over all orbits of $\times 3$ on
$\mathbb{Z}/2^j\mathbb{Z}$ for $j \leq 20$, and representative orbits
for $j \leq 30$;
(2)~direct matrix power iteration on the $2^j \times 2^j$ transfer
matrix for $j \leq 10$, tracking the geometric mean decay over complete
orbit cycles of length $L = \ord_{2^j}(3) = 2^{j-2}$.

\subsection{Phase~2c: Joint constraint analysis (\texttt{furstenberg\_lift.c})}

For each scale~$j$, we construct the combined $T_2 + T_3$ constraint
matrix (a dense $9N \times 9N$ system for the 2-scale case, with
$N = 2^j$) and compute its rank via Gaussian elimination with column
partial pivoting and tolerance $10^{-10}$.
After row reduction, we project the null space onto subsets of variables
(coarse sub-pieces, atom masses) and compute the projected null space
dimension via a second rank computation.
At $j = 10$, the matrix requires ${\sim}680$\,MB and the elimination
takes ${\sim}10$\,s on a single core.
For the 3-scale case (scales $j$, $j+1$, $j+2$), the system has
$21N$ variables and $25N$ constraints; this is limited to $j \leq 8$
(${\sim}275$\,MB).


%% ====================================================================
\section{Phase~1 Results: Orbit Structure at Scale $10^{10}$}
%% ====================================================================

\subsection{Full generation density}

We scanned $455{,}052{,}509$ primes coprime to~6 up to $10^{10}$.
Table~\ref{tab:fullgen} shows the full generation percentage by decade.

\begin{table}[h]
\centering
\caption{Full generation density $|\langle 2,3 \rangle| / (p-1) = 1$ by decade.}
\label{tab:fullgen}
\begin{tabular}{lrrr}
\toprule
Range & Primes & Full gen & Pct \\
\midrule
$[5, 10^2)$          &            23 &            18 & 78.26\% \\
$[10^2, 10^3)$       &           143 &           105 & 73.43\% \\
$[10^3, 10^4)$       &         1,061 &           756 & 71.25\% \\
$[10^4, 10^5)$       &         8,363 &         5,854 & 70.00\% \\
$[10^5, 10^6)$       &        68,906 &        48,075 & 69.77\% \\
$[10^6, 10^7)$       &       586,081 &       408,790 & 69.75\% \\
$[10^7, 10^8)$       &     5,096,876 &     3,555,045 & 69.75\% \\
$[10^8, 10^9)$       &    45,086,079 &    31,448,732 & 69.75\% \\
$[10^9, 10^{10})$    &   404,204,977 &   281,933,832 & \textbf{69.75\%} \\
\midrule
\textbf{Total}       &   455,052,509 &   317,401,207 & 69.75\% \\
\bottomrule
\end{tabular}
\end{table}

By Chebotarev's density theorem, the theoretical density of primes where
$\langle 2, 3 \rangle = (\mathbb{Z}/p\mathbb{Z})^*$ is given by the
Euler product:
\begin{equation}
\label{eq:euler}
\delta = \prod_{\ell\ \text{prime}} \left(1 - \frac{1}{\ell^2(\ell - 1)}\right)
= 0.69750136\ldots
\end{equation}
The product converges rapidly: $\ell = 2$ contributes $3/4$;
$\ell = 3$ contributes $17/18$; $\ell = 5$ contributes $99/100$.
Our empirical $69.750\%$ matches to five significant figures.

\textbf{Interpretation.}
At ``scale~$p$'' (working on $\mathbb{Z}/p\mathbb{Z}$), the
$\langle \times 2, \times 3 \rangle$ semigroup covers all of
$(\mathbb{Z}/p\mathbb{Z})^*$ for $\sim 70\%$ of primes.
For the remaining $\sim 30\%$, non-trivial invariant subsets exist,
but these are \emph{finite orbits}---precisely what Furstenberg allows.
The conjecture is consistent at every finite prime scale.

\subsection{Index distribution}

For non-full-generation primes, the index $(p-1)/|\langle 2, 3 \rangle|$
characterizes the failure mode.
Table~\ref{tab:index} shows the distribution over all $455$M primes.

\begin{table}[h]
\centering
\caption{Index distribution over $4.55 \times 10^8$ primes.}
\label{tab:index}
\begin{tabular}{rrrl}
\toprule
Index & Count & Pct & Interpretation \\
\midrule
1 & 317,401,207 & 69.75\% & Full generation \\
2 &  93,355,317 & 20.52\% & Both are QR mod $p$ \\
3 &  17,982,555 &  3.95\% & Both are cubic residues \\
4 &   9,334,173 &  2.05\% & 4th power residues \\
5 &   3,180,503 &  0.70\% & \\
6 &   4,494,649 &  0.99\% & Composite: $2 \times 3$ \\
8 &   2,722,949 &  0.60\% & \\
12 &  1,123,071 &  0.25\% & Highly composite \\
24 &    327,576 &  0.07\% & Highly composite \\
\bottomrule
\end{tabular}
\end{table}

Highly composite indices (6, 12, 24) are over-represented
because they have many prime-power divisors, each contributing
a factor to the probability.

\subsection{Maximum index growth}

Table~\ref{tab:maxindex} shows the maximum index per decade.

\begin{table}[h]
\centering
\caption{Maximum index $\max_p (p-1)/|\langle 2,3 \rangle|$ by decade.}
\label{tab:maxindex}
\begin{tabular}{lr}
\toprule
Decade & Max index \\
\midrule
$[10^2, 10^3)$ & 10 \\
$[10^3, 10^4)$ & 56 \\
$[10^4, 10^5)$ & 72 \\
$[10^5, 10^6)$ & 170 \\
$[10^6, 10^7)$ & 519 \\
$[10^7, 10^8)$ & 2,834 \\
$[10^8, 10^9)$ & 3,235 \\
$[10^9, 10^{10})$ & \textbf{35,590} \\
\bottomrule
\end{tabular}
\end{table}

The worst-case primes have very smooth $p - 1$ (many small prime factors),
causing both 2 and 3 to land in tiny subgroups.
The minimum fraction $|\langle 2, 3\rangle|/(p-1)$ in the last decade
is $2.81 \times 10^{-5}$, meaning the generators cover only
$0.003\%$ of $(\mathbb{Z}/p\mathbb{Z})^*$.

\subsection{CF convergents of $\log_2 3$}

The continued fraction $\log_2 3 = [1; 1, 1, 2, 2, 3, 1, 5, 2, 23, \ldots]$
gives the best rational approximations $p_n/q_n$ to $\log_2 3$,
equivalently the closest near-equalities $2^{p_n} \approx 3^{q_n}$.
Table~\ref{tab:cf} shows selected convergents with their cellError
and the ratio to Baker's lower bound.

\begin{table}[h]
\centering
\caption{CF convergents of $\log_2 3$ and comparison to Baker's bound.}
\label{tab:cf}
\begin{tabular}{rrrcc}
\toprule
$n$ & $a_n$ & $p_n/q_n$ & $|\operatorname{cellError}|$ & err/Baker \\
\midrule
4  & 2  & $19/12$              & $1.96 \times 10^{-2}$  & $6.6 \times 10^3$ \\
8  & 2  & $1054/665$           & $6.30 \times 10^{-5}$  & $1.8 \times 10^{10}$ \\
13 & 1  & $301994/190537$      & $9.31 \times 10^{-8}$  & $1.1 \times 10^{20}$ \\
19 & 1  & $630138897/397573379$& $1.16 \times 10^{-10}$ & $1.4 \times 10^{22}$ \\
\bottomrule
\end{tabular}
\end{table}

The err/Baker ratio grows rapidly: Baker's lower bound is exponentially
weaker than the actual CF approximation rate.
This ``gap'' is the quantitative room in which the entropy bridge must operate.

\subsection{Multi-scale cellError resonance}

On the $N \times N$ torus with threshold $\delta = 1/N$,
we count cells $(a,b)$ with $|\operatorname{cellError}(a,b)| < \delta$.
Table~\ref{tab:resonant} shows the results.

\begin{table}[h]
\centering
\caption{CellError resonant fraction at scales $N = 6^k$.}
\label{tab:resonant}
\begin{tabular}{rrrr}
\toprule
$k$ & $N = 6^k$ & Resonant cells & Fraction \\
\midrule
1 & 6         & 1 & 0.167 \\
2 & 36        & 2 & 0.056 \\
3 & 216       & 2 & 0.009 \\
4 & 1,296     & 2 & 0.0015 \\
5 & 7,776     & 3 & 0.00039 \\
6 & 46,656    & 1 & 0.000021 \\
8 & 1,679,616 & 6 & 0.0000036 \\
\bottomrule
\end{tabular}
\end{table}

The resonant count stays $O(1)$ at every scale, consistent with
Baker separation (at most $\sim 2$ CF convergents in range at any scale).
The resonant fraction decays as $\sim 1/N$, confirming that the torus
becomes overwhelmingly ``safe'' at large scales.


%% ====================================================================
\section{Phase~2 Results: The Entropy Bridge}
%% ====================================================================

\subsection{Multiplicative order asymmetry}

\begin{proposition}[Verified computationally to exponent 40]
\label{prop:orders}
For all $j \geq 3$:
\[
\ord_{2^j}(3) = 2^{j-2}, \qquad
\text{index } [\varphi(2^j) : \ord_{2^j}(3)] = 2.
\]
For all $k \geq 1$:
\[
\ord_{3^k}(2) = \varphi(3^k) = 2 \cdot 3^{k-1}, \qquad
\text{index } = 1 \;\; (\text{2 is a primitive root mod } 3^k).
\]
\end{proposition}

Both statements are classically known
(the second is Artin's conjecture restricted to $g = 2$, $p = 3$,
verified for all prime powers $3^k$).
We verify them computationally through 128-bit modular exponentiation
for all $j, k \leq 40$ (80 tests, all matching predictions).

This asymmetry is the structural foundation of the entropy bridge:
\begin{itemize}
\item $T_3$ (multiplication by~3) acts on $\mathcal{P}_2^{(j)}$ with
  period $2^{j-2}$, generating only \emph{half} of $(\mathbb{Z}/2^j\mathbb{Z})^*$.
\item $T_2$ (multiplication by~2) acts on $\mathcal{P}_3^{(k)}$ with
  full period $\varphi(3^k)$, generating \emph{all} of $(\mathbb{Z}/3^k\mathbb{Z})^*$.
\end{itemize}
The index-2 constraint on the 2-adic side means that $T_3$-invariance
forces a measure to be constant on orbits of size $2^{j-2}$,
but leaves one binary degree of freedom (the relative weight of the
two cosets).

\subsection{Joint partition atom analysis}

For balanced scales $k(j) = \lfloor j \log 2 / \log 3 + 1/2 \rfloor$
(where $2^j \approx 3^k$), the joint partition
$\mathcal{P}_2^{(j)} \vee \mathcal{P}_3^{(k)}$ has
$N = 2^j + 3^k - 1$ atoms.
Table~\ref{tab:entropy} shows the entropy statistics.

\begin{table}[h]
\centering
\caption{Entropy of $\mathcal{P}_2^{(j)} \vee \mathcal{P}_3^{(k)}$
  under Lebesgue measure (balanced scales).}
\label{tab:entropy}
\begin{tabular}{rrrcccc}
\toprule
$j$ & $k$ & Atoms & $H$ (nats) & $H/j$ & Deficit & Efficiency \\
\midrule
 4 &  3 &         42 &  3.590 & 0.897 & 0.148 & 0.960 \\
 8 &  5 &        498 &  6.020 & 0.752 & 0.191 & 0.969 \\
12 &  8 &     10,656 &  9.101 & 0.758 & 0.173 & 0.981 \\
16 & 10 &    124,584 & 11.541 & 0.721 & 0.192 & 0.984 \\
20 & 13 &  2,642,898 & 14.611 & 0.731 & 0.177 & 0.988 \\
24 & 15 & 31,126,122 & 17.063 & 0.711 & 0.190 & 0.989 \\
\bottomrule
\end{tabular}
\end{table}

\textbf{Key observation: the entropy deficit is bounded.}
Across all tested scales, the deficit $\ln(N) - H$ stays in the
interval $[0.148, 0.192]$ nats.
The efficiency $H / \ln(N)$ increases monotonically toward~1.
Under Lebesgue measure, $H/j$ converges to $\log 2 \approx 0.693$
from above (the excess comes from the additional information in the
3-adic refinement).

\begin{proposition}[Bounded entropy deficit, empirical]
\label{prop:deficit}
For all tested $(j, k)$ with $j \leq 24$ and $k = k(j)$ balanced:
\[
\ln(2^j + 3^k - 1) - H_\lambda(\mathcal{P}_2^{(j)} \vee \mathcal{P}_3^{(k)})
\leq 0.20 \text{ nats},
\]
where $\lambda$ denotes Lebesgue measure.
\end{proposition}

This bounded deficit means the atom widths are approximately uniform
at all scales.  The few thin atoms (width $\sim 1/(2^j \cdot 3^k)$,
achieving the B\'ezout floor) contribute negligibly to entropy.

\subsection{Gap anatomy}

The atoms of $\mathcal{P}_2^{(j)} \vee \mathcal{P}_3^{(k)}$ arise from
interleaving two arithmetic progressions in $[0, D)$ with $D = 2^j \cdot 3^k$:
the 2-adic boundaries $\{m \cdot 3^k\}$ (step $3^k$) and the 3-adic
boundaries $\{m' \cdot 2^j\}$ (step $2^j$).

\begin{proposition}[B\'ezout floor]
\label{prop:bezout}
For all $(j, k)$ tested (up to $j = 24$), the minimum gap between
consecutive boundary points in the common-denominator representation is
exactly~$1$.
That is,
\[
\min_{\text{atoms}} \operatorname{width} = \frac{1}{2^j \cdot 3^k}.
\]
This minimum is achieved by the B\'ezout identity: there exist
$m_0, m_0'$ with $m_0 \cdot 3^k - m_0' \cdot 2^j = 1$.
\end{proposition}

Three types of gaps occur between consecutive boundary points:

\begin{table}[h]
\centering
\caption{Gap types for $\mathcal{P}_2^{(12)} \vee \mathcal{P}_3^{(8)}$
  ($D = 2^{12} \cdot 3^8 = 26{,}873{,}856$, $N = 10{,}656$ atoms).}
\label{tab:gaps}
\begin{tabular}{lcrl}
\toprule
Gap type & Count & Width & Meaning \\
\midrule
3-adic $\to$ 3-adic & 2,464 & $2^{12} = 4096$ & 3-adic boundaries well-separated \\
Mixed (2$\leftrightarrow$3) & 8,192 & 1--4096 & Baker-governed \\
2-adic $\to$ 2-adic & 0 & --- & Impossible ($3^8 > 2^{12}$) \\
\bottomrule
\end{tabular}
\end{table}

The zero count of pure 2-adic gaps when $3^k > 2^j$ means that
every pair of 2-adic boundaries has a 3-adic boundary between them---the
3-adic partition is ``finer'' at this scale.
The mixed gaps, governed by Baker's coprimality floor, have
mean $\sim 3^k / 2$ and minimum exactly~1.

\subsection{Minimum atom width decay}

The ratio of minimum to mean atom width decays exponentially:

\begin{table}[h]
\centering
\caption{Minimum-to-mean atom width ratio (balanced scales).}
\label{tab:minmean}
\begin{tabular}{rcc}
\toprule
$j$ & min/mean & $\approx$ \\
\midrule
 8 & $8.0 \times 10^{-4}$ & $\sim 2^{-10}$ \\
12 & $4.0 \times 10^{-4}$ & $\sim 2^{-11.3}$ \\
16 & $3.2 \times 10^{-5}$ & $\sim 2^{-14.9}$ \\
20 & $1.6 \times 10^{-6}$ & $\sim 2^{-19.3}$ \\
24 & $1.3 \times 10^{-8}$ & $\sim 2^{-26.2}$ \\
\bottomrule
\end{tabular}
\end{table}

Despite the exponential shrinkage of the thinnest atom, the entropy
deficit remains bounded because there is \emph{exactly one}
B\'ezout-thin atom per partition, and its entropy contribution
$w \ln(1/w) \to 0$ as $w \to 0$.


\subsection{Transfer matrix spectral gap}

The transfer matrix of $T_3\colon x \mapsto 3x$ on the partition
$\mathcal{P}_2^{(j)}$ is the $2^j \times 2^j$ doubly stochastic matrix
$\mathbf{T}$ with entries:
\[
T_{mi} = \begin{cases} 1/3 & \text{if } m \in \{3i, 3i+1, 3i+2\} \pmod{2^j}, \\
0 & \text{otherwise.}
\end{cases}
\]
The column sums are~1 (each cell is hit by exactly 3~preimages from one source),
and the row sums are~1 (each cell distributes equally to 3~targets).

\begin{theorem}[Spectral gap of $T_3$ on $\mathcal{P}_2^{(j)}$]
\label{thm:spectral}
For $j \geq 3$, all non-trivial eigenvalues of $\mathbf{T}$ have
modulus exactly $1/3$.  The spectral gap is $1 - 1/3 = 2/3$,
independent of~$j$.
\end{theorem}

\begin{proof}[Proof sketch]
In the Fourier basis $v_k(m) = \omega^{mk}$ with $\omega = e^{2\pi i/2^j}$,
the matrix $\mathbf{T}$ maps $v_k \mapsto \lambda(k) \cdot v_{ks}$ where
$s = 3^{-1} \bmod 2^j$ and
\[
\lambda(k) = \tfrac{1}{3}\bigl(1 + 2\cos(2\pi ks / 2^j)\bigr).
\]
For an orbit $\{k_0, 3k_0, 9k_0, \ldots\}$ of size~$L$ under $\times 3 \pmod{2^j}$,
the eigenvalue modulus on the associated $L$-dimensional block is
$\bigl(\prod_{i=0}^{L-1} |\lambda(k_i)|\bigr)^{1/L}$.

Using the identity $1 + 2\cos\theta = (e^{3i\theta} - 1)/(e^{i\theta} - 1)$,
the product telescopes:
\[
\prod_{i=0}^{L-1} |1 + 2\cos(2\pi k_i s / 2^j)|
= \frac{|e^{2\pi i \cdot 3^L k_0 s / 2^j} - 1|}{|e^{2\pi i \cdot k_0 s / 2^j} - 1|}
= 1,
\]
since $3^L \equiv 1 \pmod{2^j / \gcd(k_0, 2^j)}$.
Therefore $|\text{eigenvalue}| = (1/3^L)^{1/L} = 1/3$.
\end{proof}

\textbf{Numerical verification.}
We verify this theorem in two independent ways:

\begin{enumerate}
\item \emph{Fourier orbit products} ($j = 3$ to $30$):
  For each $j \leq 20$, we enumerate \emph{all} orbits of $\times 3$ on
  $\mathbb{Z}/2^j\mathbb{Z}$ and compute $\prod |1 + 2\cos\theta_i|$;
  every orbit gives product~$= 1.0000000000$ (to 10 decimal places).
  For $j = 21$ to $30$, we check representative orbits of $k = 1$ and
  $k = 2$; the product stays within $10^{-7}$ of~1 (double-precision limit
  for sums over $10^8$ terms).

\item \emph{Direct matrix multiplication} ($j = 3$ to $10$):
  We build the full transfer matrix, apply it to an initial vector
  orthogonal to uniform, and track the geometric mean of the per-step
  decay ratio over complete orbit cycles of length $L = \text{ord}_{2^j}(3)$.
  Every cycle gives geometric mean $= 0.3333333333$ to 10 decimal places.
\end{enumerate}

\textbf{Structural observations.}
\begin{itemize}
\item The number of orbits is $2j + 1$ (for $j \leq 20$, verified),
  reflecting the $j + 1$ powers of~2 dividing $\gcd(k, 2^j)$ plus
  the coset structure.
\item The maximum orbit size is $2^{j-2} = \text{ord}_{2^j}(3)$, confirming
  Proposition~\ref{prop:orders}.
\item The eigenvalue $\lambda(2^{j-1}) = -1/3$ (the unique fixed point
  $k = 2^{j-1}$ under $\times 3$, since $3 \cdot 2^{j-1} \equiv 2^{j-1} \pmod{2^j}$).
\end{itemize}

\begin{proposition}[Exponential mixing]
\label{prop:mixing}
For any probability distribution $p$ on $\mathcal{P}_2^{(j)}$ atoms:
\[
\|\mathbf{T}^n p - \text{uniform}\|_2 \leq (1/3)^n \cdot \|p - \text{uniform}\|_2.
\]
After $n$ applications of $T_3$, the distribution is exponentially close
to uniform, at rate $1/3$ per step.
\end{proposition}

This means $T_3$-invariance at the measure level ($\mathbf{T} p = p$)
forces $p = \text{uniform}$ on $\mathcal{P}_2^{(j)}$ atoms, which gives
$H_\mu(\mathcal{P}_2^{(j)}) = j \log 2$.
The remaining question is whether $T_3$-invariance of a measure~$\mu$
on $\mathbb{T}$ implies $T_3$-invariance of the induced distribution on
$\mathcal{P}_2^{(j)}$ atoms---i.e., whether the ``lift'' from partition
to measure preserves the spectral gap constraint.

\subsection{The lift problem: straddle analysis}

The preimage $T_3^{-1}(I_m)$ of atom $I_m = [m/2^j, (m+1)/2^j)$
consists of three intervals of width $1/(3 \cdot 2^j)$:
\[
T_3^{-1}(I_m) = \bigcup_{r=0}^{2}
\left[\frac{m + r \cdot 2^j}{3 \cdot 2^j},\;
      \frac{m + r \cdot 2^j + 1}{3 \cdot 2^j}\right).
\]
In the common denominator $D = 3 \cdot 2^j$, each preimage piece
occupies the interval $[a, a+1)$ where $a = m + r \cdot 2^j$.
Since atoms in~$D$-coordinates span $[3k, 3k+3)$ and
$a + 1 \leq 3k + 3$ (as $a \bmod 3 \leq 2$),
every piece lies entirely inside one atom.
No straddling occurs: the ``straddle matrix'' $\mathbf{S}$ equals
the transfer matrix $\mathbf{T}$ exactly.

\begin{proposition}[Degrees of freedom]
\label{prop:dof}
At each scale~$j$, $T_3$-invariance of~$\mu$ imposes $N-1$
independent linear constraints on the $3N$ sub-piece masses
$\{w_{k,s}\}$ (where $k$ indexes atoms and $s \in \{0,1,2\}$ indexes
sub-pieces).
Together with the $N$ atom-sum constraints
$w_{k,0} + w_{k,1} + w_{k,2} = p_k$,
the system has rank $2N - 1$, leaving $N + 1$ degrees of freedom
(or~2 per atom plus normalization).
\end{proposition}

This is verified computationally for $j = 3$ to $12$: the constraint
rank is always $N-1$, giving exactly 2~DOF per atom at every scale.

\textbf{Interpretation.}
$T_3$-invariance at a single scale~$j$ constrains the atom masses~$p_m$
but \emph{not} the within-atom distribution.
Each atom has 3~sub-pieces (one for each preimage branch), and the
relative weights of these sub-pieces are free parameters.
The transfer matrix equation $\mathbf{T} p = p$ is the \emph{special case}
where each sub-piece carries exactly $1/3$ of its atom's mass.

\textbf{Multi-scale constraints.}
At scale~$j+1$, each atom of~$\mathcal{P}_2^{(j)}$ splits into two
atoms of~$\mathcal{P}_2^{(j+1)}$.  Each old sub-piece (width $1/(3 \cdot 2^j)$)
spans $\sim 4/3$ new atoms, so the scale-$(j+1)$ constraints
partially fix the within-atom distribution at scale~$j$.
After $\lceil \log_2 3 \rceil + 1 = 3$ refinements, each old sub-piece
is resolved into separate atoms, and the scale-$(j+3)$ spectral gap
forces the resolved masses toward uniformity.

The question is: do \emph{multi-scale} constraints (T$_3$-invariance
at scales $j$ and $j+1$ simultaneously, plus $T_2$-invariance
linking the scales) eliminate the within-atom DOFs?
This is the subject of the next section.


\subsection{Joint $T_2 + T_3$ constraint rank analysis}
\label{sec:lift}

We now build the full constraint system for joint $T_2$- and
$T_3$-invariance across multiple dyadic scales, and compute its rank
to determine the degrees of freedom exactly.

\subsubsection*{Mathematical setup (2-scale system)}

Consider two consecutive scales~$j$ and $j+1$, with $N = 2^j$.
The sub-piece masses are:
\begin{align*}
\text{Scale } j\!: &\quad w_{k,s} \;\text{ for }
  k \in [0, N),\; s \in \{0,1,2\}
  &&\text{($3N$ variables)}, \\
\text{Scale } j\!+\!1\!: &\quad w'_{m,s} \;\text{ for }
  m \in [0, 2N),\; s \in \{0,1,2\}
  &&\text{($6N$ variables)}.
\end{align*}
The total is $9N$ unknowns.
Four blocks of homogeneous linear constraints relate them:

\begin{enumerate}
\item \textbf{$T_3$-invariance at scale~$j$} ($N$ equations):
  for each atom~$m$,
  \[
  \sum_{s=0}^{2} w_{m,s}
  = \sum_{r=0}^{2} w_{\lfloor(m+rN)/3\rfloor,\;(m+rN)\bmod 3}.
  \]

\item \textbf{$T_3$-invariance at scale~$j\!+\!1$} ($2N$ equations):
  same structure with $2N$ atoms and $6N$ fine variables.

\item \textbf{Refinement} ($3N$ equations):
  each coarse sub-piece is the union of two fine sub-pieces.
  Explicitly:
  \begin{align*}
  w_{k,0} &= w'_{2k,0} + w'_{2k,1}, \\
  w_{k,1} &= w'_{2k,2} + w'_{2k+1,0}, \\
  w_{k,2} &= w'_{2k+1,1} + w'_{2k+1,2}.
  \end{align*}
  These follow from the geometric subdivision: sub-piece $(k,s)$
  at scale~$j$ covers $[(3k+s)/(3 \cdot 2^j),\; (3k+s+1)/(3 \cdot 2^j))$,
  which at scale~$j+1$ (in coordinates with denominator $3 \cdot 2^{j+1}$)
  spans $[6k+2s,\; 6k+2s+2)$ and splits at $6k+2s+1$.

\item \textbf{$T_2$-linking} ($3N$ equations):
  $T_2$-invariance relates sub-pieces across scales via
  \[
  w_{k,s} = w'_{k,s} + w'_{k+N,s}
  \qquad (0 \le k < N,\; s \in \{0,1,2\}),
  \]
  since $T_2^{-1}$ of a sub-piece at scale~$j$ decomposes into
  two sub-pieces at scale~$j+1$ separated by~$N$.
\end{enumerate}

\subsubsection*{Rank computation}

We compute the rank of the constraint matrix via Gaussian elimination.
Table~\ref{tab:lift2} shows the results for $j = 3$ to~$10$.

\begin{table}[h]
\centering
\caption{2-scale constraint rank analysis ($9N$ variables,
  verified for $j = 3, \ldots, 10$).}
\label{tab:lift2}
\begin{tabular}{rrrrrrr}
\toprule
$j$ & $N$ & $9N$ & DOF (no $T_2$) & DOF ($+T_2$) & $T_2$ kills & Atom DOF \\
\midrule
 3 &      8 &     72 &    33 &    17 &    16 &    5 \\
 4 &     16 &    144 &    65 &    33 &    32 &    9 \\
 5 &     32 &    288 &   129 &    65 &    64 &   17 \\
 6 &     64 &    576 &   257 &   129 &   128 &   33 \\
 7 &    128 &  1{,}152 &   513 &   257 &   256 &   65 \\
 8 &    256 &  2{,}304 & 1{,}025 &   513 &   512 &  129 \\
 9 &    512 &  4{,}608 & 2{,}049 & 1{,}025 & 1{,}024 &  257 \\
10 &  1{,}024 &  9{,}216 & 4{,}097 & 2{,}049 & 2{,}048 &  513 \\
\bottomrule
\end{tabular}
\end{table}

\begin{proposition}[Joint constraint rank, 2-scale]
\label{prop:liftrank}
For the 2-scale system at scales $j$, $j+1$ with $N = 2^j$ ($j \ge 3$):
\begin{alignat*}{2}
\operatorname{rank}(\text{no } T_2) &= 5N - 1,
  &\qquad \operatorname{DOF}(\text{no } T_2) &= 4N + 1, \\
\operatorname{rank}(\text{with } T_2) &= 7N - 1,
  &\qquad \operatorname{DOF}(\text{with } T_2) &= 2N + 1.
\end{alignat*}
$T_2$-invariance kills exactly $2N$ degrees of freedom.
\end{proposition}

\subsubsection*{3-scale extension}

The same analysis extends to three consecutive scales ($j$, $j+1$, $j+2$)
with $21N$ variables and $25N$ constraint rows.
Table~\ref{tab:lift3} shows the results.

\begin{table}[h]
\centering
\caption{3-scale constraint rank analysis ($21N$ variables,
  verified for $j = 3, \ldots, 8$).}
\label{tab:lift3}
\begin{tabular}{rrrrrr}
\toprule
$j$ & $N$ & $21N$ & DOF (no $T_2$) & DOF ($+T_2$) & $T_2$ kills \\
\midrule
3 &    8 &    168 &    65 &    33 &    32 \\
4 &   16 &    336 &   129 &    65 &    64 \\
5 &   32 &    672 &   257 &   129 &   128 \\
6 &   64 &  1{,}344 &   513 &   257 &   256 \\
7 &  128 &  2{,}688 & 1{,}025 &   513 &   512 \\
8 &  256 &  5{,}376 & 2{,}049 & 1{,}025 & 1{,}024 \\
\bottomrule
\end{tabular}
\end{table}

The 3-scale rank formulas are:
\begin{alignat*}{2}
\operatorname{rank}(\text{no } T_2) &= 13N - 1,
  &\qquad \operatorname{DOF}(\text{no } T_2) &= 8N + 1, \\
\operatorname{rank}(\text{with } T_2) &= 17N - 1,
  &\qquad \operatorname{DOF}(\text{with } T_2) &= 4N + 1.
\end{alignat*}
In both cases, the total DOF follows the universal formula
\begin{equation}
\label{eq:dof}
\operatorname{DOF}(\text{with } T_2)
= N_{\text{finest}} + 1,
\end{equation}
where $N_{\text{finest}}$ is the number of atoms at the finest scale
($2N$ for 2-scale, $4N$ for 3-scale).
The extra DOF beyond~1 (overall scaling) resides at the finest scale,
which has no child constraints.

\subsubsection*{Projection analysis}

To understand where the freedom lives, we project the null space
of the full 2-scale system (with~$T_2$) onto subsets of variables.
Table~\ref{tab:proj} shows the projected null space dimension at
each level.

\begin{table}[h]
\centering
\caption{Projected null space dimension for the 2-scale system
  (pattern exact for $j = 3, \ldots, 10$).}
\label{tab:proj}
\begin{tabular}{lcc}
\toprule
Projection & Null dim & Formula \\
\midrule
Full system ($9N$ variables) & $2N + 1$ & $N_{\text{finest}} + 1$ \\
Coarse sub-pieces ($3N$ variables) & $N + 1$ & \\
Coarse atom masses ($N$ variables) & $N/2 + 1$ & \\
\bottomrule
\end{tabular}
\end{table}

The atom-mass DOF at the coarsest scale is $N/2 + 1$,
and this value is \emph{unchanged} between the 2-scale and 3-scale
systems (verified for $j = 3, \ldots, 8$).
Adding more scales does not further constrain the coarsest-scale
atom masses within the linear system.

\subsubsection*{Interpretation}

The atom-mass DOF $= N/2 + 1 \gg 2$ means the \emph{linear} constraint
system does not force $\mathbf{T}p = p$ (which would give DOF~$= 2$,
corresponding to Lebesgue and~$\delta_0$).
The extra $N/2 - 1$ dimensions correspond to signed
``pseudo-measures'' that satisfy all the $T_3$, $T_2$, and refinement
constraints algebraically but violate positivity ($w_{k,s} \ge 0$).

This identifies \textbf{positivity} as the essential ingredient for the
entropy bridge.
The linear-algebraic structure establishes:
\begin{itemize}
\item $T_2$-invariance halves the DOF at every scale
  (killing exactly $N_{\text{finest}}$ directions);
\item The spectral gap (Theorem~\ref{thm:spectral}) forces any
  \emph{genuine} solution of $\mathbf{T}p = p$ to be uniform;
\item But the measure-level $T_3$-invariance constraint is weaker than
  $\mathbf{T}p = p$, leaving ${\sim}\,N/2$ directions of freedom
  in atom masses that are incompatible with positivity but satisfy
  the linear system.
\end{itemize}

Both Lebesgue measure (uniform: $w_{k,s} = 1/(3N)$ at scale~$j$,
$w'_{m,s} = 1/(6N)$ at scale~$j+1$) and $\delta_0$ (point mass:
$w_{0,0} = 1$, $w'_{0,0} = 1$, all others~$0$) are verified to
be exact solutions of the full constraint system.
These span a 2-dimensional subspace within the $(2N+1)$-dimensional
null space.
The remaining $2N - 1$ null directions have negative components
and therefore cannot represent probability measures.


%% ====================================================================
\section{Synthesis: Architecture of the Entropy Bridge}
%% ====================================================================

The Phase~1 and Phase~2 data reveal the following architecture for
the entropy bridge argument:

\begin{enumerate}
\item \textbf{Baker coprimality floor.}
  The minimum distance between a 2-adic and a 3-adic boundary is
  $1 / (2^j \cdot 3^k)$, achieved exactly by B\'ezout.
  Atoms never collapse.

\item \textbf{Bounded entropy deficit.}
  $\ln(N) - H \leq 0.20$ nats for all tested balanced scales up to
  $j = 24$.
  The atom widths are approximately uniform under Lebesgue measure,
  with the non-uniformity entirely explained by a single B\'ezout-thin atom.

\item \textbf{Order asymmetry.}
  $\ord_{2^j}(3) = 2^{j-2}$ (index~2), while $\ord_{3^k}(2) = \varphi(3^k)$ (index~1).
  Any $T_3$-invariant measure on $\mathcal{P}_2^{(j)}$ must be constant
  on orbits of size $2^{j-2}$, leaving only two cosets as free parameters.

\item \textbf{Full generation at prime scales.}
  For $69.75\%$ of primes~$p$, $\langle 2, 3 \rangle = (\mathbb{Z}/p\mathbb{Z})^*$,
  leaving no room for non-trivial invariant measures.
  For the remaining primes, invariant subsets exist but are
  \emph{finite orbits}---exactly what Furstenberg allows.

\item \textbf{Spectral gap $= 2/3$.}
  The transfer matrix of $T_3$ on $\mathcal{P}_2^{(j)}$ has all
  non-trivial eigenvalues with modulus exactly $1/3$
  (Theorem~\ref{thm:spectral}).
  This gives exponential mixing at rate $1/3$ per step,
  so $T_3$-invariance at the partition level forces exact uniformity.

\item \textbf{Joint constraint rank $= 7N - 1$.}
  The combined $T_2 + T_3$ constraint system at two consecutive
  scales has rank $7N - 1$ out of $9N$ variables
  (Proposition~\ref{prop:liftrank}), leaving DOF $= 2N + 1$.
  $T_2$-invariance kills exactly $2N$ directions, halving the
  freedom at every scale.
  However, the \emph{linear} system does not force the transfer
  equation $\mathbf{T}p = p$: the atom-mass DOF is $N/2 + 1$,
  not~$2$.
  The extra directions are signed pseudo-measures that violate
  positivity.
\end{enumerate}

\subsection{The gap to close}

Points (1)--(5) together show that:
\begin{itemize}
\item Baker separation ensures atoms never collapse (point~1).
\item The entropy deficit is bounded under Lebesgue measure (point~2).
\item $T_3$-invariance at the partition level forces exact uniformity on
  $\mathcal{P}_2^{(j)}$ atoms via the spectral gap (points~3, 5).
\end{itemize}

The remaining question is whether $T_3$-invariance of a \emph{measure}~$\mu$
implies $T_3$-invariance of the induced distribution on partition atoms.
Specifically, if $\mu$ is $T_3$-invariant, does the atom mass vector
$p_m = \mu([m/2^j, (m+1)/2^j))$ satisfy $\mathbf{T} p = p$?

If $\mu$ is $T_3$-invariant and the partition $\mathcal{P}_2^{(j)}$ were
$T_3$-invariant (i.e., $T_3^{-1}(\mathcal{P}_2^{(j)}) \preccurlyeq \mathcal{P}_2^{(j)}$),
then yes.  But $\mathcal{P}_2^{(j)}$ is a $T_2$-partition, not a $T_3$-partition,
so the refinement $T_3^{-1}(\mathcal{P}_2^{(j)})$ generally creates atoms
that straddle the 2-adic boundaries.

Point~(6) sharpens this obstacle.
The joint $T_2 + T_3$ constraint system, even across multiple scales,
leaves $N/2 + 1$ degrees of freedom in the atom masses---far more
than the~$2$ needed for $\mathbf{T}p = p$.
The extra directions are signed pseudo-measures that satisfy all
linear constraints but violate positivity ($w_{k,s} \ge 0$).
Adding more scales does not reduce this freedom: the atom-mass
DOF is $N/2 + 1$ for both 2-scale and 3-scale systems.

This identifies \textbf{positivity} as the essential non-linear ingredient.
The linear-algebraic framework determines the constraint architecture
(exact ranks, DOF formulas), but the final step---showing that
the only \emph{non-negative} solutions are Lebesgue and~$\delta_0$---requires
a convexity or extremal-measure argument that goes beyond rank computation.
One promising route is to exploit the spectral gap: since $\mathbf{T}$
contracts non-uniform components by factor $1/3$, any non-uniform
positive eigenvector of the multi-scale refinement must amplify
signed components at finer scales, eventually violating positivity.


%% ====================================================================
\section{Comparison with Prior Work}
%% ====================================================================

Our computational-Diophantine approach differs from the existing
literature in several respects:

\begin{itemize}
\item \textbf{Rudolph (1990)} proves the conjecture under $h_\mu > 0$
  using symbolic dynamics and Rohklin partitions.
  We aim to show that Baker's effective bounds \emph{force} $h_\mu > 0$,
  reducing to Rudolph rather than replacing it.

\item \textbf{EKL (2006)} use measure rigidity on homogeneous spaces
  to show any counterexample has dimension~0.
  Our dimension argument (Phase~3, future work) would recover this
  result by an entirely different route: multi-scale Baker separation
  on the torus.

\item \textbf{No prior work} has combined Baker's theorem with explicit
  partition entropy computation.
  The bounded entropy deficit (Proposition~\ref{prop:deficit})
  and the order asymmetry (Proposition~\ref{prop:orders}) are,
  to our knowledge, new observations in this context.
\end{itemize}


%% ====================================================================
\section{Honest Assessment}
%% ====================================================================

\textbf{Difficulty: Extreme.}
Furstenberg's conjecture has resisted all approaches for over 55~years.
The gap between ``Baker separation makes resonances sparse''
(which we have demonstrated) and ``Baker separation forces entropy
production'' (which we conjecture) is precisely the gap that
separates the known results (Rudolph, EKL) from the full conjecture.

\textbf{What we have accomplished:}
\begin{itemize}
\item Verified the Euler product prediction to five significant figures
  over $4.55 \times 10^8$ primes.
\item Demonstrated bounded entropy deficit across six orders of magnitude
  in partition scale.
\item Established the order asymmetry as a structural invariant of the
  $\times 2, \times 3$ system.
\item Proved and verified the spectral gap $= 2/3$ theorem
  (Theorem~\ref{thm:spectral}), showing exponential mixing of $T_3$
  on $\mathcal{P}_2^{(j)}$.
\item Computed exact rank formulas for the joint $T_2 + T_3$
  multi-scale constraint system (Proposition~\ref{prop:liftrank}),
  identifying positivity as the essential non-linear ingredient
  for the entropy bridge.
\item Built a computational framework (OpenMP, checkpointed, extensible
  to GPU) for further investigation.
\end{itemize}

\textbf{Realistic best case:}
We prove the Entropy Production Conjecture, reducing Furstenberg
to Rudolph.  This would be a major theorem.

\textbf{Realistic worst case:}
We produce a partial result (e.g., Furstenberg for measures
satisfying an explicit Diophantine condition on their support),
publishable independently, that clarifies the remaining obstacle.


%% ====================================================================
\begin{thebibliography}{99}

\bibitem{baker1966}
A.~Baker, \emph{Linear forms in the logarithms of algebraic numbers},
Mathematika \textbf{13} (1966), 204--216.

\bibitem{einsiedler2006}
M.~Einsiedler, A.~Katok, and E.~Lindenstrauss,
\emph{Invariant measures and the set of exceptions to Littlewood's conjecture},
Ann.\ of Math.\ \textbf{164} (2006), 513--560.

\bibitem{furstenberg1967}
H.~Furstenberg, \emph{Disjointness in ergodic theory, minimal sets,
and a problem in Diophantine approximation},
Math.\ Systems Theory \textbf{1} (1967), 1--49.

\bibitem{host1995}
B.~Host, \emph{Nombres normaux, entropie, translations},
Israel J.\ Math.\ \textbf{91} (1995), 75--83.

\bibitem{rhin1987}
G.~Rhin, \emph{Approximants de Pad\'e et mesures effectives
d'irrationalit\'e},
S\'eminaire de Th\'eorie des Nombres, Paris 1985--86,
Progr.\ Math.\ \textbf{71}, Birkh\"auser, 1987, pp.~155--164.

\bibitem{rudolph1990}
D.~J.~Rudolph, \emph{$\times 2$ and $\times 3$ invariant measures
and entropy},
Ergodic Theory Dynam.\ Systems \textbf{10} (1990), 395--406.

\end{thebibliography}

\end{document}
